\documentclass[12pt,a4paper]{article}

% ── Paquetes ──────────────────────────────────────────────────────────────────
\usepackage[utf8]{inputenc}
\usepackage[T1]{fontenc}
\usepackage[spanish]{babel}
\usepackage{geometry}
\geometry{margin=2.5cm}
\usepackage{lmodern}
\usepackage{microtype}
\usepackage{parskip}
\usepackage{setspace}
\onehalfspacing

\usepackage{listings}
\usepackage{xcolor}
\usepackage{booktabs}
\usepackage{array}
\usepackage{longtable}
\usepackage{hyperref}
\usepackage{amsmath}
\usepackage{amssymb}
\usepackage{graphicx}
\usepackage{fancyhdr}
\usepackage{titlesec}
\usepackage{enumitem}
\usepackage{mdframed}

% ── Configuración de página ───────────────────────────────────────────────────
\pagestyle{fancy}
\fancyhf{}
\rhead{Jorge Rodríguez González}
\lhead{Adaptación de MIDI2ScoreTransformer}
\cfoot{\thepage}

% ── Colores para código ───────────────────────────────────────────────────────
\definecolor{codebg}{RGB}{248,248,248}
\definecolor{codeframe}{RGB}{200,200,200}
\definecolor{kwcolor}{RGB}{0,0,200}
\definecolor{strcolor}{RGB}{180,0,0}
\definecolor{cmtcolor}{RGB}{80,120,80}
\definecolor{diffadd}{RGB}{220,255,220}
\definecolor{diffdel}{RGB}{255,220,220}

\lstdefinestyle{python}{
  language=Python,
  backgroundcolor=\color{codebg},
  frame=single,
  rulecolor=\color{codeframe},
  basicstyle=\ttfamily\small,
  keywordstyle=\color{kwcolor}\bfseries,
  stringstyle=\color{strcolor},
  commentstyle=\color{cmtcolor}\itshape,
  showstringspaces=false,
  breaklines=true,
  breakatwhitespace=true,
  tabsize=4,
  numbers=left,
  numberstyle=\tiny\color{gray},
  numbersep=8pt,
  captionpos=b,
  literate={á}{{\'a}}1 {é}{{\'e}}1 {í}{{\'i}}1 {ó}{{\'o}}1 {ú}{{\'u}}1
           {ñ}{{\~n}}1 {Á}{{\'A}}1 {É}{{\'E}}1 {Í}{{\'I}}1 {Ó}{{\'O}}1
           {Ú}{{\'U}}1 {Ñ}{{\~N}}1 {ü}{{\"u}}1 {Ü}{{\"U}}1,
}

\lstset{style=python}

% ── Entorno para notas de cambio ──────────────────────────────────────────────
\newmdenv[
  backgroundcolor=yellow!10,
  linecolor=orange!60,
  linewidth=1.2pt,
  roundcorner=4pt,
  innerleftmargin=10pt,
  innerrightmargin=10pt,
  innertopmargin=6pt,
  innerbottommargin=6pt,
]{notebox}

\newmdenv[
  backgroundcolor=red!5,
  linecolor=red!50,
  linewidth=1.2pt,
  roundcorner=4pt,
  innerleftmargin=10pt,
  innerrightmargin=10pt,
]{errorbox}

\newmdenv[
  backgroundcolor=green!5,
  linecolor=green!50,
  linewidth=1.2pt,
  roundcorner=4pt,
  innerleftmargin=10pt,
  innerrightmargin=10pt,
]{fixbox}

% ── Título ────────────────────────────────────────────────────────────────────
\title{
  \vspace{-1cm}
  \Large\textbf{Informe técnico de adaptación del módulo \texttt{roformer.py}}\\[0.4em]
  \large Proyecto: MIDI2ScoreTransformer ---
  Conversión de MIDI de interpretación a partitura\\[0.4em]
  \normalsize Actualización de compatibilidad con \texttt{transformers} v5.x
}
\author{
  Jorge Rodríguez González\\
  \small Grado en Ingeniería Informática\\
}
\date{Febrero 2026}

% ══════════════════════════════════════════════════════════════════════════════
\begin{document}
% ══════════════════════════════════════════════════════════════════════════════

\maketitle
\thispagestyle{fancy}

\tableofcontents
\newpage

% ──────────────────────────────────────────────────────────────────────────────
\section{Introducción y motivación}
% ──────────────────────────────────────────────────────────────────────────────

El presente informe documenta los cambios realizados sobre el fichero \texttt{models/roformer.py}
del repositorio \textit{MIDI2ScoreTransformer} \cite{beyer2024}, publicado por Tim Beyer y Angela Dai
como código asociado al artículo presentado en la conferencia ISMIR 2024. El objetivo original del
sistema es la conversión automática de archivos MIDI de interpretación (\textit{performance MIDI})
a partituras musicales en formato MusicXML, utilizando un modelo Transformer encoder-decoder.

Al intentar ejecutar el código en un entorno con las versiones actuales de las bibliotecas
de \textit{deep learning}, concretamente con \texttt{transformers==5.2.0} y
\texttt{lightning>=2.0.0}, el sistema falla al intentar realizar inferencia. El error concreto
que se produce es el siguiente:

\begin{errorbox}
\texttt{TypeError: CustomRoFormerSelfAttention.forward() got an unexpected keyword argument 'past\_key\_values'}
\end{errorbox}

La causa raíz es una incompatibilidad entre la arquitectura de herencia de clases utilizada
por el autor y los cambios internos introducidos en \texttt{transformers} v5.x. En este
informe se analiza en detalle el problema, se describe la solución adoptada y se justifican
todas las decisiones de diseño tomadas durante el proceso de adaptación.

% ──────────────────────────────────────────────────────────────────────────────
\section{Descripción de la arquitectura original}
% ──────────────────────────────────────────────────────────────────────────────

\subsection{Estrategia de herencia selectiva}

El fichero \texttt{roformer.py} original implementa un modelo RoFormer
\cite{su2021roformer} personalizado mediante una técnica de \textbf{herencia selectiva}
sobre las clases internas de la biblioteca HuggingFace \texttt{transformers}.
La idea es importar las clases base del RoFormer estándar y sobreescribir selectivamente
aquellos componentes que el autor desea modificar, dejando intactos los demás.

Las modificaciones específicas que el autor introduce respecto al RoFormer estándar son:
\begin{enumerate}[label=\alph*)]
  \item \textbf{Pre-norm}: la normalización de capa (\textit{LayerNorm}) se aplica
        \textit{antes} de la atención, en lugar de después (post-norm), lo que mejora
        la estabilidad del entrenamiento.
  \item \textbf{Activación SwiGLU}: la red de alimentación directa
        (\textit{Feed-Forward Network}) sustituye la activación GELU estándar por
        SwiGLU \cite{shazeer2020glu}, que combina una compuerta sigmoide con
        la función SiLU.
  \item \textbf{Fusión QKV}: las proyecciones de \textit{query}, \textit{key} y
        \textit{value} del self-attention se fusionan en una única capa lineal para
        mayor eficiencia.
  \item \textbf{Flash Attention}: se utiliza \texttt{F.scaled\_dot\_product\_attention}
        de PyTorch, que aprovecha la implementación eficiente en memoria de Flash
        Attention cuando está disponible.
  \item \textbf{RoPE en cross-attention}: las codificaciones posicionales rotatorias
        se aplican también en el módulo de cross-attention del decoder.
\end{enumerate}

\subsection{Cadena de herencia original}

La jerarquía de clases del fichero original es la siguiente:

\begin{center}
\begin{tabular}{ll}
\toprule
\textbf{Clase HuggingFace (base)} & \textbf{Clase personalizada (autor)} \\
\midrule
\texttt{nn.Module} & \texttt{CustomRoFormerSelfAttention} \\
\texttt{RoFormerAttention} & \texttt{CustomRoFormerAttention} \\
\texttt{nn.Module} & \texttt{CustomRoFormerIntermediate} \\
\texttt{nn.Module} & \texttt{CustomRoFormerOutput} \\
\texttt{RoFormerLayer} & \texttt{CustomRoFormerLayer} \\
\texttt{RoFormerEncoder} & \texttt{CustomRoFormerEncoder} \\
\texttt{RoFormerModel} (HF) & \texttt{RoFormerModel} (autor) \\
\texttt{pl.LightningModule} & \texttt{Roformer} \\
\bottomrule
\end{tabular}
\end{center}

El problema crítico reside en que \texttt{RoFormerAttention}, \texttt{RoFormerEncoder}
y \texttt{RoFormerLayer} son clases \textbf{internas} de HuggingFace, accesibles a través
de la ruta \texttt{transformers.models.roformer.modeling\_roformer}, que no forma parte
de la API pública estable de la biblioteca.

% ──────────────────────────────────────────────────────────────────────────────
\section{Análisis del error: cambio de API en \texttt{transformers} v5.x}
% ──────────────────────────────────────────────────────────────────────────────

\subsection{El interceptor \texttt{modeling\_layers.py}}

En \texttt{transformers} v5.x, HuggingFace introdujo una nueva capa de abstracción
denominada \texttt{GradientCheckpointingLayer} (implementada en \texttt{modeling\_layers.py})
que \textbf{intercepta todas las llamadas \texttt{\_\_call\_\_}} de los módulos que heredan
de sus clases base internas. Este interceptor normaliza y traduce los nombres de los
argumentos antes de pasarlos al método \texttt{forward()} del módulo.

Concretamente, el interceptor renombra el argumento \texttt{past\_key\_value}
(en singular, convención interna del código del autor) a \texttt{past\_key\_values}
(en plural, nueva convención de HuggingFace v5). Como el método
\texttt{CustomRoFormerSelfAttention.forward()} solo declara \texttt{past\_key\_value}
(singular) en su firma, Python lanza el error de argumento inesperado.

\subsection{Por qué ocurre exactamente en \texttt{CustomRoFormerLayer}}

El error se produce en la cadena de llamadas siguiente:

\begin{enumerate}
  \item \texttt{CustomRoFormerEncoder.forward()} llama a \texttt{layer\_module(...)}
        (que es una instancia de \texttt{CustomRoFormerLayer}).
  \item \texttt{CustomRoFormerLayer} \textbf{no} tiene \texttt{forward()} propio en el
        código original, por lo que se ejecuta el de su clase padre:
        \texttt{RoFormerLayer.forward()} de HuggingFace.
  \item \texttt{RoFormerLayer.forward()} de HF llama a \texttt{self.attention(...)}
        pasando \texttt{past\_key\_values} (plural).
  \item El interceptor de \texttt{modeling\_layers.py} actúa sobre esta llamada porque
        \texttt{CustomRoFormerAttention} hereda de \texttt{RoFormerAttention} (clase HF).
  \item El argumento llega a \texttt{CustomRoFormerSelfAttention.forward()} como
        \texttt{past\_key\_values} (plural), pero la firma esperada es
        \texttt{past\_key\_value} (singular) $\Rightarrow$ \texttt{TypeError}.
\end{enumerate}

El problema es, por tanto, una \textbf{combinación de herencia implícita y cambio de
convención de nombres} entre versiones mayores de \texttt{transformers}.

% ──────────────────────────────────────────────────────────────────────────────
\section{Solución adoptada: desacoplamiento completo de HuggingFace}
% ──────────────────────────────────────────────────────────────────────────────

La solución adoptada consiste en \textbf{eliminar completamente la herencia de las
clases internas de HuggingFace} y reimplementar toda la jerarquía como subclases
directas de \texttt{nn.Module} de PyTorch. Dado que el autor ya había sobreescrito
prácticamente toda la lógica relevante, esta operación no implica pérdida de
funcionalidad, sino únicamente una reorganización estructural.

Las ventajas de este enfoque son:
\begin{itemize}
  \item \textbf{Estabilidad}: el código ya no depende de internals de HuggingFace
        que pueden cambiar en cualquier versión.
  \item \textbf{Transparencia}: toda la lógica del modelo reside en un único fichero,
        sin comportamientos heredados implícitamente.
  \item \textbf{Compatibilidad de pesos}: los nombres de las capas en el
        \texttt{state\_dict} no cambian, por lo que el checkpoint preentrenado
        carga sin modificaciones.
\end{itemize}

% ──────────────────────────────────────────────────────────────────────────────
\section{Descripción detallada de cada cambio}
% ──────────────────────────────────────────────────────────────────────────────

\subsection{Cambio 1: Eliminación de imports de internals de HuggingFace}

\textbf{Código original:}
\begin{lstlisting}[caption={Imports originales en roformer.py}]
from transformers import RoFormerModel as RoFormerModelBase
from transformers.modeling_outputs import \
    BaseModelOutputWithPastAndCrossAttentions
from transformers.models.roformer.modeling_roformer import (
    RoFormerAttention,
    RoFormerEncoder,
    RoFormerLayer
)
from models.embedding import MIDIEmbeddings, MXLEmbeddings, MXLUnembeddings
from models.model import BaseModel
\end{lstlisting}

\textbf{Código nuevo:}
\begin{lstlisting}[caption={Imports corregidos en roformer.py}]
from transformers import RoFormerConfig
from transformers.modeling_outputs import \
    BaseModelOutputWithPastAndCrossAttentions
import numpy as np
# Resto: solo PyTorch estandar (torch, torch.nn, torch.nn.functional)
\end{lstlisting}

\begin{notebox}
\textbf{Justificación:} \texttt{RoFormerAttention}, \texttt{RoFormerEncoder} y
\texttt{RoFormerLayer} pertenecen a la API interna de HuggingFace, no a su API pública.
Su comportamiento puede cambiar sin aviso entre versiones. Se mantiene
\texttt{BaseModelOutputWithPastAndCrossAttentions} porque es un dataclass de la API
pública estable que simplifica el retorno de los outputs del encoder. \texttt{RoFormerConfig}
también se mantiene porque el checkpoint lo usa como clase de configuración serializada.
\end{notebox}

\subsection{Cambio 2: Implementación propia de \texttt{RoFormerSinusoidalPositionalEmbedding}}

En el código original esta clase no existía en el fichero: se instanciaba implícitamente
dentro del \texttt{\_\_init\_\_} de \texttt{RoFormerEncoder} de HuggingFace al llamar a
\texttt{super().\_\_init\_\_(config)}. Al eliminar esa herencia, es necesario implementarla
explícitamente.

\begin{lstlisting}[caption={Implementación propia del embedding posicional sinusoidal}]
class RoFormerSinusoidalPositionalEmbedding(nn.Embedding):
    def __init__(self, num_positions: int, embedding_dim: int,
                 padding_idx=None):
        super().__init__(num_positions, embedding_dim)
        self._init_weight()

    def _init_weight(self):
        n_pos, dim = self.weight.shape
        position_enc = np.array([
            [pos / np.power(10000, 2 * (j // 2) / dim)
             for j in range(dim)]
            for pos in range(n_pos)
        ])
        out = torch.empty(n_pos, dim, dtype=self.weight.dtype,
                          requires_grad=False)
        sentinel = dim // 2 if dim % 2 == 0 else (dim // 2) + 1
        out[:, 0:sentinel] = torch.FloatTensor(
            np.sin(position_enc[:, 0::2]))
        out[:, sentinel:]  = torch.FloatTensor(
            np.cos(position_enc[:, 1::2]))
        self.weight = nn.Parameter(out, requires_grad=False)

    @torch.no_grad()
    def forward(self, input_ids_shape, past_key_values_length: int = 0,
                position_ids=None):
        if position_ids is None:
            bsz, seq_len = input_ids_shape[:2]
            position_ids = torch.arange(
                past_key_values_length,
                past_key_values_length + seq_len,
                dtype=torch.long,
                device=self.weight.device,
            )
        return super().forward(position_ids)
\end{lstlisting}

\begin{notebox}
\textbf{Justificación matemática:} Las codificaciones posicionales rotatorias (RoPE)
se calculan mediante frecuencias sinusoidales:
\[
  \text{PE}_{(p,2i)}   = \sin\!\left(\frac{p}{10000^{2i/d}}\right), \quad
  \text{PE}_{(p,2i+1)} = \cos\!\left(\frac{p}{10000^{2i/d}}\right)
\]
donde $p$ es la posición en la secuencia y $d$ la dimensión del embedding por cabeza
de atención ($d = \texttt{hidden\_size} / \texttt{num\_attention\_heads} = 64$).
El parámetro \texttt{sentinel} es el índice que separa la mitad de senos de la mitad
de cosenos en el tensor de pesos. El atributo \texttt{requires\_grad=False} es
deliberado: estos embeddings son fijos y no se aprenden durante el entrenamiento.
\end{notebox}

\subsection{Cambio 3: \texttt{CustomRoFormerSelfAttention} --- sin cambios en lógica}

Esta clase ya era \texttt{nn.Module} puro en el código original, por lo que su lógica
interna se mantiene íntegramente. El cambio relevante es estructural: al no existir
ninguna clase padre de HuggingFace en la jerarquía, el interceptor
\texttt{modeling\_layers.py} de \texttt{transformers} v5.x \textbf{no actúa} sobre
esta clase, y el argumento \texttt{past\_key\_value} (singular) fluye correctamente
desde el \texttt{forward()} hasta el KV-cache.

\begin{notebox}
\textbf{Punto clave:} El nombre \texttt{past\_key\_value} (en singular) es una
convención interna del código del autor, no un error tipográfico. HuggingFace adoptó
\texttt{past\_key\_values} (plural) en v4.x como nombre estándar. El interceptor de
v5.x impone este estándar a todos los módulos que hereden de sus clases, pero al ser
\texttt{CustomRoFormerSelfAttention} un \texttt{nn.Module} puro, esta imposición no
aplica.
\end{notebox}

\subsection{Cambio 4: \texttt{CustomRoFormerAttention} --- de herencia a composición}

\textbf{Código original:}
\begin{lstlisting}[caption={CustomRoFormerAttention original (herencia de HF)}]
class CustomRoFormerAttention(RoFormerAttention):
    def __init__(self, config, is_cross_attention=False):
        super().__init__(config)  # crea capas de HF que luego se descartan
        self.self = CustomRoFormerSelfAttention(config, is_cross_attention)
        self.output = CustomRoFormerSelfOutput(config)
\end{lstlisting}

\textbf{Código nuevo:}
\begin{lstlisting}[caption={CustomRoFormerAttention corregida (composicion pura)}]
class CustomRoFormerAttention(nn.Module):
    def __init__(self, config, is_cross_attention=False):
        super().__init__()
        self.self = CustomRoFormerSelfAttention(config, is_cross_attention)
        self.output = CustomRoFormerSelfOutput(config)
        self.pruned_heads = set()

    def forward(self, hidden_states, attention_mask, sinusoidal_pos,
                head_mask, encoder_hidden_states, encoder_attention_mask,
                past_key_value, output_attentions):
        self_outputs = self.self(
            hidden_states, attention_mask, sinusoidal_pos, head_mask,
            encoder_hidden_states, encoder_attention_mask,
            past_key_value, output_attentions,
        )
        attention_output = self.output(self_outputs[0], hidden_states)
        outputs = (attention_output,) + self_outputs[1:]
        return outputs
\end{lstlisting}

\begin{notebox}
\textbf{Justificación:} La llamada \texttt{super().\_\_init\_\_(config)} de
\texttt{RoFormerAttention} instanciaba capas de self-attention y output propias de HF
que se sobreescribían inmediatamente en las líneas siguientes. Era código muerto que
consumía memoria innecesariamente. El atributo \texttt{self.pruned\_heads = set()} se
añade por completitud, ya que algunos mecanismos de serialización de modelos pueden
esperarlo. La adición del método \texttt{forward()} explícito elimina la dependencia
del \texttt{forward()} de la clase padre de HF.
\end{notebox}

\subsection{Cambio 5: \texttt{CustomRoFormerLayer} --- adición de \texttt{forward()} propio}

Este es el \textbf{cambio más crítico} de toda la adaptación. La clase original heredaba
de \texttt{RoFormerLayer} pero \textbf{no declaraba un método \texttt{forward()} propio},
lo que significa que en tiempo de ejecución se utilizaba el \texttt{forward()} de
\texttt{RoFormerLayer} de HuggingFace, el cual pasa \texttt{past\_key\_values} (plural)
al llamar a \texttt{self.attention()}.

\textbf{Código original:}
\begin{lstlisting}[caption={CustomRoFormerLayer original (sin forward propio)}]
class CustomRoFormerLayer(RoFormerLayer):
    def __init__(self, config):
        super().__init__(config)
        self.attention = CustomRoFormerAttention(config)
        if self.add_cross_attention:
            self.crossattention = CustomRoFormerAttention(
                config, is_cross_attention=True)
        self.intermediate = CustomRoFormerIntermediate(config)
        self.output = CustomRoFormerOutput(config)
    # No tiene forward() --> usa el de RoFormerLayer de HF --> BUG
\end{lstlisting}

\textbf{Código nuevo:}
\begin{lstlisting}[caption={CustomRoFormerLayer corregida (con forward completo)}]
class CustomRoFormerLayer(nn.Module):
    def __init__(self, config):
        super().__init__()
        self.attention = CustomRoFormerAttention(config)
        self.is_decoder = config.is_decoder
        self.add_cross_attention = config.add_cross_attention
        if self.add_cross_attention:
            self.crossattention = CustomRoFormerAttention(
                config, is_cross_attention=True)
        self.intermediate = CustomRoFormerIntermediate(config)
        self.output = CustomRoFormerOutput(config)

    def forward(self, hidden_states, attention_mask, sinusoidal_pos,
                head_mask, encoder_hidden_states, encoder_attention_mask,
                past_key_value=None, output_attentions=False):

        # Separar KV-cache: [:2] para self-attn, [-2:] para cross-attn
        self_attn_pkv = past_key_value[:2] \
            if past_key_value is not None else None
        self_attention_outputs = self.attention(
            hidden_states, attention_mask, sinusoidal_pos, head_mask,
            output_attentions=output_attentions,
            past_key_value=self_attn_pkv,
        )
        attention_output = self_attention_outputs[0]

        if self.is_decoder:
            outputs = self_attention_outputs[1:-1]
            present_key_value = self_attention_outputs[-1]
        else:
            outputs = self_attention_outputs[1:]

        if self.is_decoder and encoder_hidden_states is not None:
            cross_attn_pkv = past_key_value[-2:] \
                if past_key_value is not None else None
            cross_attention_outputs = self.crossattention(
                attention_output, attention_mask, sinusoidal_pos,
                head_mask, encoder_hidden_states, encoder_attention_mask,
                cross_attn_pkv, output_attentions,
            )
            attention_output = cross_attention_outputs[0]
            outputs = outputs + cross_attention_outputs[1:-1]
            present_key_value = present_key_value \
                + cross_attention_outputs[-1]

        intermediate_output = self.intermediate(attention_output)
        layer_output = self.output(intermediate_output, attention_output)
        outputs = (layer_output,) + outputs
        if self.is_decoder:
            outputs = outputs + (present_key_value,)
        return outputs
\end{lstlisting}

\begin{notebox}
\textbf{Justificación del particionado del KV-cache:} En un decoder Transformer con
cross-attention, el KV-cache de una capa almacena cuatro tensores: los dos primeros
corresponden a las claves y valores del \textit{self-attention} causal, y los dos
últimos a las claves y valores del \textit{cross-attention} con el encoder. La
indexación \texttt{[:2]} y \texttt{[-2:]} distribuye correctamente estos tensores a
cada submódulo de atención.
\end{notebox}

\subsection{Cambio 6: \texttt{CustomRoFormerEncoder} --- herencia de \texttt{nn.Module}}

\textbf{Código original:}
\begin{lstlisting}[caption={CustomRoFormerEncoder original}]
class CustomRoFormerEncoder(RoFormerEncoder):
    def __init__(self, config):
        super().__init__(config)  # instancia embed_positions de HF
        self.layer = nn.ModuleList(
            [CustomRoFormerLayer(config)
             for _ in range(config.num_hidden_layers)])
        # el forward() si era propio del autor (se mantiene intacto)
\end{lstlisting}

\textbf{Código nuevo:}
\begin{lstlisting}[caption={CustomRoFormerEncoder corregida}]
class CustomRoFormerEncoder(nn.Module):
    def __init__(self, config):
        super().__init__()
        self.config = config
        self.embed_positions = RoFormerSinusoidalPositionalEmbedding(
            config.max_position_embeddings,
            config.hidden_size // config.num_attention_heads,
        )
        self.layer = nn.ModuleList(
            [CustomRoFormerLayer(config)
             for _ in range(config.num_hidden_layers)])
        self.gradient_checkpointing = False
\end{lstlisting}

\begin{notebox}
\textbf{Justificación:} Al dejar de heredar de \texttt{RoFormerEncoder}, es necesario
instanciar explícitamente \texttt{embed\_positions}. La dimensión del embedding
posicional es \texttt{hidden\_size // num\_attention\_heads} = $512 / 8 = 64$, porque
RoPE opera sobre los vectores de query y key de cada cabeza de atención
individualmente, no sobre el embedding completo. El \texttt{forward()} de esta clase
era ya propio del autor en el código original y se conserva sin cambios.
\end{notebox}

\subsection{Cambio 7: \texttt{RoFormerModel} --- de \texttt{PreTrainedModel} a \texttt{nn.Module}}

\textbf{Código original:}
\begin{lstlisting}[caption={RoFormerModel original (hereda de PreTrainedModel de HF)}]
class RoFormerModel(RoFormerModelBase):  # RoFormerModelBase = PreTrainedModel
    def __init__(self, config):
        super().__init__(config)
        self.encoder = CustomRoFormerEncoder(config)
        del self.embeddings  # elimina embeddings de HF no necesarios
\end{lstlisting}

\textbf{Código nuevo:}
\begin{lstlisting}[caption={RoFormerModel corregida con metodos propios}]
class RoFormerModel(nn.Module):
    def __init__(self, config):
        super().__init__()
        self.config = config
        self.encoder = CustomRoFormerEncoder(config)

    def get_extended_attention_mask(self, attention_mask, input_shape):
        if attention_mask.dim() == 2:
            if self.config.is_decoder:
                B, T = input_shape
                seq_ids = torch.arange(T, device=attention_mask.device)
                causal_mask = (seq_ids[None, None, :]
                               .repeat(B, T, 1)
                               <= seq_ids[None, :, None])
                causal_mask = causal_mask[:, None, :, :].to(
                    attention_mask.dtype)
                extended = (causal_mask
                            * attention_mask[:, None, None, :])
                extended = ((1.0 - extended)
                            * torch.finfo(torch.float32).min)
            else:
                extended = ((1.0 - attention_mask[:, None, None, :])
                            * torch.finfo(torch.float32).min)
        else:
            raise ValueError(
                f"attention_mask debe ser 2D, es {attention_mask.dim()}D")
        return extended

    def invert_attention_mask(self, encoder_attention_mask):
        if encoder_attention_mask.dim() == 2:
            enc_ext = encoder_attention_mask[:, None, None, :]
        else:
            raise ValueError(
                f"encoder_attention_mask debe ser 2D")
        enc_ext = ((1.0 - enc_ext.float())
                   * torch.finfo(torch.float32).min)
        return enc_ext
\end{lstlisting}

\begin{notebox}
\textbf{Justificación:} \texttt{PreTrainedModel} de HuggingFace es una clase con
cientos de métodos auxiliares (\texttt{from\_pretrained}, \texttt{save\_pretrained},
\texttt{push\_to\_hub}, etc.) de los cuales el autor solo necesitaba dos:
\texttt{get\_extended\_attention\_mask} e \texttt{invert\_attention\_mask}. Ambos se
reimplementan directamente.

\texttt{get\_extended\_attention\_mask} convierte la máscara binaria $(B, T)$, donde
$1$ indica token real y $0$ indica padding, en una máscara $(B, 1, T, T)$ compatible
con la operación de suma sobre las puntuaciones de atención, donde las posiciones
enmascaradas toman el valor $-\infty$ (implementado como el mínimo del tipo
\texttt{float32}). Para el decoder, se combina con una máscara causal triangular
inferior que impide atender a posiciones futuras.

\texttt{invert\_attention\_mask} realiza la misma conversión pero para la máscara del
encoder en el cross-attention del decoder.
\end{notebox}

\subsection{Cambio 8: \texttt{forward\_enc} y \texttt{forward\_dec} --- parámetros hardcodeados}

\textbf{Código original:}
\begin{lstlisting}[caption={Lectura de parametros desde el config (version original)}]
encoder_outputs = self.encoder.encoder(
    embedding_output,
    ...
    output_attentions=self.enc_config.output_attentions,
    output_hidden_states=self.enc_config.output_hidden_states,
    return_dict=self.enc_config.use_return_dict,
)
\end{lstlisting}

\textbf{Código nuevo:}
\begin{lstlisting}[caption={Parametros hardcodeados (version corregida)}]
encoder_outputs = self.encoder.encoder(
    embedding_output,
    ...
    output_attentions=False,
    output_hidden_states=False,
    return_dict=True,
)
\end{lstlisting}

\begin{notebox}
\textbf{Justificación:} En \texttt{transformers} v5.x, \texttt{RoFormerConfig}
convirtió \texttt{output\_attentions} de un atributo simple a una \textit{property}
de Python que internamente accede a \texttt{self.\_output\_attentions}. Al reconstruir
el objeto \texttt{MyModelConfig} desde el checkpoint guardado con
\texttt{transformers==4.29.2}, el atributo privado \texttt{\_output\_attentions} no
existe porque la versión antigua no lo usaba, produciéndose el error
\texttt{AttributeError: 'MyModelConfig' object has no attribute '\_output\_attentions'}.

Durante la inferencia, estos tres valores son constantes: no se necesitan los mapas de
atención (\texttt{output\_attentions=False}), no se necesitan los estados intermedios
(\texttt{output\_hidden\_states=False}), y siempre se usa el formato de retorno en
diccionario (\texttt{return\_dict=True}). Hardcodearlos es, por tanto, la solución
correcta, y además mejora ligeramente el rendimiento al evitar ramas de código
condicionales.
\end{notebox}

% ──────────────────────────────────────────────────────────────────────────────
\section{Verificación: compatibilidad con el checkpoint preentrenado}
% ──────────────────────────────────────────────────────────────────────────────

Un requisito fundamental de la adaptación es que los pesos del modelo preentrenado
(fichero \texttt{MIDI2ScoreTF.ckpt}, guardado con PyTorch Lightning y
\texttt{transformers==4.29.2}) carguen correctamente en la nueva arquitectura.

Las primeras y últimas claves del \texttt{state\_dict} del checkpoint son:

\begin{center}
\small
\begin{tabular}{l}
\toprule
\textbf{Claves del state\_dict (checkpoint preentrenado)} \\
\midrule
\texttt{encoder.encoder.embed\_positions.weight} \\
\texttt{encoder.encoder.layer.0.attention.self.query\_key\_value.weight} \\
\texttt{encoder.encoder.layer.0.attention.self.query\_key\_value.bias} \\
\texttt{encoder.encoder.layer.0.attention.self.norm.weight} \\
\texttt{encoder.encoder.layer.0.attention.self.norm.bias} \\
\ldots \\
\texttt{embeddings\_enc.embeddings.unconditional.weight} \\
\texttt{embeddings\_enc.embeddings.velocity.weight} \\
\texttt{embeddings\_enc.embeddings.velocity.bias} \\
\texttt{embeddings\_enc.layer\_norm.weight} \\
\texttt{embeddings\_enc.layer\_norm.bias} \\
\bottomrule
\end{tabular}
\end{center}

Estas claves coinciden exactamente con la estructura de la arquitectura modificada,
ya que los \textbf{nombres de todas las capas se han conservado} durante la
refactorización. La llamada \texttt{model.load\_state\_dict(ck["state\_dict"],
strict=True)} no produce ninguna clave faltante ni inesperada.

El modelo resultante tiene \textbf{32.595.514 parámetros} y se ejecuta correctamente
tanto en CPU como en GPU (CUDA), confirmado con la siguiente salida de inferencia sobre
un MIDI real de 1172 notas:

\begin{fixbox}
\texttt{Cargando checkpoint: MIDI2ScoreTF.ckpt}\\
\texttt{Modelo listo en cuda | Params: 32,595,514}\\
\texttt{Procesando MIDI: prelude\_1.mid}\\
\texttt{Notas en el MIDI: 1172}\\
\texttt{Generando... 100\% (1172/1172 notas)}\\
\texttt{Tokens generados: 1172 notas}\\
\texttt{Convirtiendo tokens a MusicXML...}\\
\texttt{\checkmark\ Partitura guardada en: prelude\_1.musicxml}
\end{fixbox}

% ──────────────────────────────────────────────────────────────────────────────
\section{Resumen de cambios}
% ──────────────────────────────────────────────────────────────────────────────

\begin{longtable}{>{\ttfamily\small}p{5.2cm}
                  >{\small}p{3.5cm}
                  >{\small}p{5.8cm}}
\toprule
\textbf{Clase / método} & \textbf{Cambio} & \textbf{Motivo} \\
\midrule
\endhead
Imports de internals HF & Eliminados & Inestabilidad entre versiones mayores \\
\addlinespace
RoFormerSinusoidal\-Positional\-Embedding & Nueva clase & Instanciación
  implícita en HF eliminada \\
\addlinespace
CustomRoFormerSelf\-Attention & Sin cambios lógicos & Ya era nn.Module puro;
  el bug venía de la jerarquía superior \\
\addlinespace
CustomRoFormerAttention & Herencia $\to$ composición & super().\_\_init\_\_
  creaba capas HF inmediatamente descartadas \\
\addlinespace
CustomRoFormerLayer & Adición de forward() & Sin forward propio, usaba el de
  HF que pasaba past\_key\_values (plural) \\
\addlinespace
CustomRoFormerEncoder & Herencia $\to$ nn.Module & Necesario instanciar
  embed\_positions manualmente \\
\addlinespace
RoFormerModel & PreTrainedModel $\to$ nn.Module & Solo se necesitaban 2
  métodos de PreTrainedModel; se reimplementan \\
\addlinespace
forward\_enc / forward\_dec & Parámetros hardcodeados & output\_attentions
  deviene property en v5.x; \_output\_attentions no existe en checkpoints
  antiguos \\
\bottomrule
\caption{Resumen de todos los cambios realizados en roformer.py}
\end{longtable}

% ──────────────────────────────────────────────────────────────────────────────
\section{Conclusiones}
% ──────────────────────────────────────────────────────────────────────────────

La adaptación realizada resuelve la incompatibilidad del sistema
\textit{MIDI2ScoreTransformer} con \texttt{transformers} v5.x mediante una estrategia
de \textbf{desacoplamiento estructural}: se elimina la dependencia de las clases
internas de HuggingFace y se convierte toda la jerarquía del modelo en módulos PyTorch
puros. Los cambios son quirúrgicos en el sentido de que no modifican ningún peso,
ninguna operación matemática ni ningún hiperparámetro del modelo original; únicamente
reorganizan la estructura de herencia en Python.

El resultado es un sistema que:
\begin{itemize}
  \item Carga correctamente el checkpoint preentrenado sin remapar ninguna clave.
  \item Ejecuta inferencia completa sobre MIDIs reales produciendo partituras en
        formato MusicXML.
  \item Es estable frente a futuras actualizaciones de \texttt{transformers}, al no
        depender de ninguna API interna de la biblioteca.
  \item Mantiene todas las mejoras arquitectónicas del autor: pre-norm, SwiGLU,
        fusión QKV, Flash Attention y RoPE en cross-attention.
\end{itemize}

% ──────────────────────────────────────────────────────────────────────────────
\begin{thebibliography}{9}
\bibitem{beyer2024}
  T.~Beyer y A.~Dai,
  \textit{End-to-End Piano Performance-MIDI to Score Conversion with Transformers},
  Proc.\ 25th Int.\ Society for Music Information Retrieval Conf.\ (ISMIR),
  San Francisco, 2024.

\bibitem{su2021roformer}
  J.~Su, Y.~Lu, S.~Pan, A.~Murtadha, B.~Wen y Y.~Liu,
  \textit{RoFormer: Enhanced Transformer with Rotary Position Embedding},
  arXiv:2104.09864, 2021.

\bibitem{shazeer2020glu}
  N.~Shazeer,
  \textit{GLU Variants Improve Transformer},
  arXiv:2002.05202, 2020.

\bibitem{vaswani2017}
  A.~Vaswani et al.,
  \textit{Attention Is All You Need},
  Advances in Neural Information Processing Systems, vol.~30, 2017.
\end{thebibliography}

\end{document}

